\PassOptionsToPackage{dvipsnames}{color}
\PassOptionsToPackage{dvipsnames}{xcolor}
\documentclass[10pt]{NSP1}
\usepackage{url,floatflt}
\usepackage{helvet,times}
\usepackage{psfig,graphics}
\usepackage{mathptmx,amsmath,amssymb,bm}
\usepackage{float}
\usepackage[bf,hypcap]{caption}
%% --- additional packages required by the tables and figures of this paper ---
\usepackage{graphicx}
\usepackage{longtable,booktabs,array,calc}
\usepackage{multirow}

\providecommand{\tightlist}{\setlength{\itemsep}{0pt}\setlength{\parskip}{0pt}}
\providecommand{\real}[1]{#1}
\providecommand{\ul}[1]{\underline{#1}}
\providecommand{\hl}[1]{#1}
\providecommand{\RL}[1]{#1}
\providecommand{\textquotesingle}{'}
\DeclareUnicodeCharacter{03B1}{$\alpha$}
\DeclareUnicodeCharacter{03B2}{$\beta$}
\DeclareUnicodeCharacter{03C7}{$\chi$}
\DeclareUnicodeCharacter{03C1}{$\rho$}
\DeclareUnicodeCharacter{03BC}{$\mu$}
\DeclareUnicodeCharacter{03C3}{$\sigma$}
\DeclareUnicodeCharacter{2264}{$\leq$}
\DeclareUnicodeCharacter{2265}{$\geq$}
\DeclareUnicodeCharacter{2260}{$\neq$}
\DeclareUnicodeCharacter{00D7}{$\times$}
\DeclareUnicodeCharacter{00B1}{$\pm$}
\DeclareUnicodeCharacter{2248}{$\approx$}
\DeclareUnicodeCharacter{221A}{$\sqrt{\ }$}
\DeclareUnicodeCharacter{2192}{$\rightarrow$}
\DeclareUnicodeCharacter{2022}{$\bullet$}
\DeclareUnicodeCharacter{00B7}{$\cdot$}
\DeclareUnicodeCharacter{2032}{$'$}
\DeclareUnicodeCharacter{0393}{$\Gamma$}
\DeclareUnicodeCharacter{03B7}{$\eta$}
\DeclareUnicodeCharacter{03B8}{$\theta$}
\DeclareUnicodeCharacter{03BB}{$\lambda$}
\DeclareUnicodeCharacter{03A3}{$\Sigma$}
\DeclareUnicodeCharacter{03B4}{$\delta$}
\DeclareUnicodeCharacter{03B3}{$\gamma$}
\DeclareUnicodeCharacter{2026}{\ldots}
\DeclareUnicodeCharacter{FEFF}{}

\tolerance=1
\emergencystretch=\maxdimen
\hyphenpenalty=10000
\hbadness=10000

\topmargin=0.00cm

\def\sm{\smallskip}
\def\no{\noindent}

\def\firstpage{1}
\setcounter{page}{\firstpage}
\def\thevol{--}
\def\thenumber{--}
\def\theyear{20--}

\begin{document}

\titlefigurecaption{{\large \bf \rm Journal of Statistics Applications \& Probability }\\ {\it\small An International Journal}}

\title{The Role of Value-Centered Counseling in Enhancing Decision-Making Skills among Jordanian University Students}

\author{Maram Jaser Bani Salameh$^{1}$}
\institute{$^{1}$Department of Psychological Counseling, Faculty of Educational Sciences, Jadara University, Irbid, Jordan, Dr.maram077@gmail.com}

\titlerunning{Value-centered counseling and decision-making skills}
\authorrunning{M. J. Bani Salameh}

%corresponding author email
\mail{Dr.maram077@gmail.com}

\received{7 Jun. 2026}
\revised{20 Jun. 2026}
\accepted{25 Jul. 2026}
\published{1 Aug. 2026}

\abstracttext{The purpose of this research was to examine the effect of value-centered counseling on improving decision-making skills among students in Jordanian universities. This study used the descriptive and correlation quantitative research methodology. Data was gathered using a questionnaire survey that targeted 384 Jordanian university students. Value-centered counseling was operational zed into the following dimensions: value clarification, self-awareness, personal responsibility, social and moral commitment, and future orientation. Skills related to decision-making were measured using problem identification, gathering information, evaluating alternative solutions, choosing the right decision and being confident and responsible in making decisions. It was found that students had high levels of value-centered counseling and decision-making skills. There were statistically significant positive correlations between all dimensions of value-centered counseling and decision-making skills. It was also found that value-centered counseling had a significant regression effect on decision-making skills, where value clarification, personal responsibility and future orientation were the most significant predictors. Based on the findings of this research, value-centered counseling can improve the students' skills in making value-based, responsible and future-oriented decisions.}

\keywords{Value-centered counseling, decision-making skills, university students, values clarification, personal responsibility, Jordanian universities}

\maketitle

\raggedbottom
\setlength{\textfloatsep}{12pt plus 2pt minus 2pt}
\setlength{\intextsep}{12pt plus 2pt minus 2pt}
\setlength{\floatsep}{12pt plus 2pt minus 2pt}

\section{Introduction}

University students face different academic, personal, social and future career choices. Hence, decision-making is an essential life skill that will help the students cope with challenges at university, find alternatives and take responsibility for their actions \cite{ref1}. In addition, recent studies have found that decision-making is closely related to self-awareness, responsibility, social-emotional competence and capability to predict the outcomes before making decisions \cite{ref2}.

The value-centered counseling is an important counseling technique as it helps students understand their personal values, discover what is important to them and make decisions based on their beliefs and goals. Value clarification contributes to building the confidence and readiness for decision-making, while counseling based on values develops social and moral qualities such as empathy and responsibility \cite{ref3,ref4}.

In this respect, the current study explores the impact of value-centered counseling on development of decision-making skills in Jordanian university students.

\subsection{Research problem}

Even though university students are expected to be able to make mature decisions, some students might hesitate, lack self-efficacy, depend on others or have problems with evaluating alternatives. The previous literature proves that maladaptive decision-making style leads to procrastination among university students \cite{ref5,ref6}. While counseling and guidance help students' career decision-making and self-efficacy \cite{ref7}.

So, the research question is stated in the following way:

What is the role of value-centered counseling in improving decision-making skills among Jordanian university students?

\subsection{Research questions}

1. What is the level of value-centered counseling among Jordanian university students?

2. What is the level of decision-making skills among Jordanian university students?

3. Is there any statistically significant correlation between value-centered counseling and decision-making skills?

4. Does value-centered counseling have any statistically significant influence on decision-making skills?

5. What dimensions of value-centered counseling contribute the most to development of decision-making skills?

\subsection{Research objectives}

The purpose of this research is to:

1. Determining the level of value-centered counseling among Jordanian university students.

2. Determining the level of decision-making skills among Jordanian university students.

3. Investigating the correlation between value-centered counseling and decision-making skills.

4. Determining the effect of value-centered counseling on decision-making skills.

5. Finding out the most influential dimensions of value-centered counseling for developing students' decision-making skills.

\subsection{Research hypotheses}

H1: There is a statistically significant positive correlation between value-centered counseling and decision-making skills among Jordanian university students.

H2: Value-centered counseling has a statistically significant effect on decision-making skills among Jordanian university students.

H2a: Value clarification has a statistically significant effect on decision-making skills.

H2b: Self-awareness has a statistically significant effect on decision-making skills.

H2c: Personal responsibility has a statistically significant effect on decision-making skills.

H2d: Moral and social commitment has a statistically significant effect on decision-making skills.

H2e: Future orientation has a statistically significant effect on decision-making skills.

\subsection{Significance of the study}

The current study is important in its focus on university students, who are in the crucial period of their academic, personal and career development. The results of the study may provide the information useful for university counseling practice and design of programs that would develop students' values, responsibility, self-awareness and future orientation. Besides, the study can also provide practical criteria for Jordanian universities to improve students' decision-making skills using counseling services. Recent studies highlight the importance of values in students' decisions and career preferences of young people \cite{ref8,ref9}.

\section{Theoretical Framework}

\subsection{Value-Centered Counseling}

Value-centered counseling assumes that the maturity of students' choices depends on their articulation of personal values, self-understanding, and responsible behavior. This approach does not imply value imposition but rather helps the students to articulate their beliefs, find out what matters to them and figure out how these values could guide their decisions at university \cite{ref10}. Values are inner standards which define behavior and make decisions meaningful in this model. Current literature on counseling and psychology states that working with values enables the transition from unclear reaction to meaningful and conscious action \cite{ref11}. Value-centered counseling also correlates with self-awareness and behavior because the students need to know themselves in order to make decisions in line with their goals and identity \cite{ref12}. Value-centered counseling is especially relevant in the university setting where students experience identity formation, adjustment to the educational institution and future plans. During this stage, counseling is supposed to assist not only in solving current problems but also to provide the students with a stable inner reference point for judgments and decisions \cite{ref13}. Self-determination theory holds that autonomous behavior appears as a result of people perceiving their actions as personally meaningful rather than forced \cite{ref14}. Value-centered counseling could thus be understood as a process that helps students to become more autonomous, responsible and act in accordance with consciously chosen values.

As a result, this research defines value-centered counseling through five dimensions. Value clarification is the student's ability to define values and priorities which guide his or her behavior. Self-awareness is the student's understanding of his or her motivations, strengths, weaknesses and needs. Personal responsibility is the student's ability to accept consequences of his or her choices. Moral and social commitment is the student's consideration of ethical standards and rights of other people when making decisions. Future orientation is the student's ability to link current decisions with future academic, personal and professional goals. These are the five dimensions which constitute the theoretical basis of value-centered counseling used in this study.

\subsection{Decision-Making Skills}

Decision-making skills imply the student's ability to overcome uncertainty and make a rational choice through understanding the situation, evaluation of information, comparison of alternatives and taking responsibility for the final decision. Decision-making is not a single event but a cognitive and personal process guided by information, values, self-confidence and the ability to assess consequences. \cite{ref15} Gati and Kulcsár characterize successful decision-making as finding promising alternatives, their evaluation, comparison and the selection of the most fitting option, thus seeing decision-making as a structured process not only a reaction under the influence of some external force \cite{ref16}.

University students need decision-making skills because of the number of decisions to be made regarding study specialization, academic performance, relationships, time management and future career. Inefficiency of decision-making skills could be demonstrated by hesitation \cite{ref17}. dependence on others, poor evaluation of alternatives, fear of responsibility and socially motivated decisions. The recent theoretical research on career decision-making shows that problems with decision-making usually appear due to the lack of readiness or information which is either insufficient or contradictory \cite{ref17}. Showing that good decision-making implies both readiness of the person and information. The same is true about higher education students who need self-understanding, information analysis and self-confidence to be prepared to make significant decisions \cite{ref18}.

In this study, decision-making skills are characterized by five dimensions. Problem identification is the recognition and definition of the decision situation. Information gathering is the student's ability to gather relevant and reliable information prior to selecting. Evaluation of alternatives is the comparison of possible options according to their advantages, disadvantages and consequences \cite{ref19}. Choosing the appropriate decision is the ability to choose the option which fits the situation and the student's values. Confidence and responsibility in decision-making is the student's ability to trust his or her decision and accept its consequences. These are the five dimensions which constitute the dependent variable in the current research \cite{ref20}.

The theoretical relationship between the two variables is based on the idea that values guide the process of decision-making. When the student clarifies his or her values, he or she is better able to distinguish between alternatives; when he or she develops self-awareness, he or she understands motives of his or her decisions; when he or she increases his or her responsibility and moral commitment, he or she is more careful with the evaluation of consequences; finally, future orientation helps him or her to take into account long-term consequences and to avoid impulsive decisions.

Thus, this study suggests that value-centered counseling could help to increase decision-making skills of Jordanian university students providing them with the inner frame of judgment. The relationship can be expressed as follows: students demonstrating higher levels of value clarification, self-awareness, responsibility, moral and social commitment and future orientation are expected to demonstrate better skills in identifying problems, gathering information, evaluating alternatives, choosing the appropriate decision and accepting responsibility for his or her choices.

\subsection{Previous Studies}

Literature review of recent years on value-centered counseling and decision-making may be viewed through two interrelated trajectories. The first one refers to the values clarification, value-based counseling, and the role of values in facilitating students' self-awareness, commitment, and action. The second trajectory focuses on decision-making skills, especially among university students, viewing decision-making as a process conditioned by confidence, self-efficacy, information, values, and future orientation \cite{ref21}.

First of all, there is increasing amount of research devoted to values not only as individuals' beliefs but as guidelines to actions. In case of university setting, the fact is especially important since the students always make various decisions regarding academic and professional matters while constructing their identity at the same time. Specifically, \cite{ref22} Bayly and Bumpus have studied values clarity among first-year college students and regarded it as an important aspect of adaptation and personal growth. Thus, values clarification is not just a question of morality but an important practical concept which helps the students to understand their actions.

More precisely, in terms of counseling, the study of \cite{ref23} Viskovich, Pakenham, and Fowler has analyzed experiential exercises related to values and committed action among university students. The researchers have found out that such exercises could attract the attention of the students and stimulate them to think about their performance-related goals and priorities and committed action. It is important to note that the present research implies that value-centered counseling will influence the process of decision-making not by providing certain prescriptions but by giving students an internal mechanism of making responsible choices.

Moreover, Larsson, Hartley, and McHugh \cite{ref24} have conducted a research concerning the effect of a brief web-based Acceptance and Commitment Therapy (ACT) intervention on the college students' psychological state. Despite the fact that the study was mostly dedicated to the problem of mental health, it is interesting for the purpose of the present paper since the ACT-based interventions involve value, acceptance, and committed action. Thus, the idea that value-based counseling may be used among the university students as a way of development and prevention is not limited to treating psychological problems.

Finally, Hsu, Adamowicz, and Thomas \cite{ref25} have provided more evidence in the form of systematic review and meta-analysis of the effect of ACT interventions among undergraduate students. Their findings have shown the positive effect on psychological flexibility and inflexibility. The fact is important since the psychological flexibility is closely connected with the student's ability to act in line with his/her values regardless of pressure, uncertainty, and emotional stress. In this sense, value-centered counseling may help students to make better decisions by increasing their tolerance to uncertainty and encouraging them to take actions in accordance with their values.

There were several studies which focused on the connection between the values clarification and professional and ethical decision-making. For example, De Gagne \cite{ref26}. Has analyzed the effect of values clarification exercises among nursing students who prepare for integration of artificial intelligence. The author has pointed out that values clarification could help the students to reflect upon their ethical responsibilities and choices in professional life. Though the situation differs from the counseling of university students, it confirms the same idea: values clarification assists in moving from passivity to judgment.

More directly, Heo \cite{ref27} has conducted a research concerning the effect of career value clarification program on career consciousness maturity and career decision-making among female college students. The findings of the researcher showed that value clarification could lead to more mature career awareness and more rational career decision-making. The finding is very relevant to the present research since it shows how value clarification affects decision-making among college students.

The research on decision-making skills also proves that students' decisions are influenced not only by the information. In particular, Wang, Wang, and Lai \cite{ref28} have conducted the analysis of sustainable career development among college students from the perspective of the Social Cognitive Career Theory. The findings showed that career values, career decision-making self-efficacy, and career goals were important for shaping students' career decision-making. It is important for the purpose of the present research since it confirms the idea that values could affect the decision-making process indirectly by increasing confidence, orientation, and future orientation.

Wang, Wang, and WenYa \cite{ref28} have also conducted the analysis of the quality of career decision-making among the students of Chinese higher vocational colleges. The authors have found out that career values and career decision-making self-efficacy had the positive effect on the quality of students' career decisions. The finding is significant since it means that the quality of decisions improves if both internal factors (values and self-efficacy) are developed simultaneously. Thus, decision-making process should be viewed not just as a technical skill but as value- and confidence-based.

Also, Lee, Jung, Baek, and Lee \cite{ref1} have examined the connection between career decision-making self-efficacy, career preparation behavior, and career decision difficulties among South Korean college students. The researchers have found out that the students with high level of decision-making self-efficacy had more career preparation behavior and less career decision-making difficulties. The finding is relevant to the present research since it means that confidence and responsibility should be included among the dimensions of decision-making skills since the students lacking confidence would know the alternative actions but fail to make decision.

Similar results have been provided by Liu, Zhang, Dang, and Gao \cite{ref29}. They have found out that career education skills were connected with career adaptability, career decision-making self-efficacy being a mediator. It adds an important dimension to the analysis: decision-making is not just a final choice but the part of the adaptive process. The students need to be educated, have self-efficacy and adaptability to navigate the changing academic and professional conditions.

The similar idea has been confirmed by the study of Salim, Istiasih, Rumalutur, and Situmorang \cite{ref30} . These researchers have found out that career decision self-efficacy mediated the connection between the peer support and career adaptability among final-year college students. It means that the external support would not be enough unless the students believe in their decision-making skills. It is relevant to the value-centered counseling in the sense that counseling should not just provide advice but help the students to develop confidence, responsibility, and orientation. The study of Zhou, Liu, Xu, and Jobe \cite{ref31} has also confirmed the impact of the social support on career decision-making difficulties and the fact that the psychological capital and career decision-making self-efficacy have mediating role. The research supports the present study's idea that decision-making difficulties are connected with psychological resources like hope, confidence, resilience, and self-efficacy.

Pham, Lam, and Bui \cite{ref32} have also demonstrated that career exploration has moderated the connection between self-efficacy and career choice among university students. The study emphasizes the importance of the exploration and social support in order to turn the confidence into choices. It reinforces the current framework in the sense that value-centered counseling could help students to explore different alternatives based on meaningful criteria and not under pressure or imitation.

The similar conclusion has been drawn by Milledzi, Asamani, Opoku Mensah, and Naamwanuru \cite{ref33}. They have confirmed the mediation role of the career decision self-efficacy in the connection between the general self-efficacy and career decision-making. The finding means that the students could have general confidence but still lack specific decision-making confidence in order to make right choices. It is relevant to the decision-making dimension of the present study, especially to the confidence and responsibility in decision-making.

In the Arab and Jordanian context, Daheri \cite{ref34}. has conducted the analysis of the effectiveness of a group counseling program in developing decision-making skills among students in the Counseling and Mental Health Department at World Islamic Sciences and Education University. \cite{ref38} The research has shown the statistically significant improvement in the experimental group after the program implementation. Though the study is relevant to the topic, it addresses the group counseling in general and not value-centered counseling as a multidimensional variable. Thus, the present research becomes necessary.

Summing up, the previous research has shown that values, self-efficacy, psychological readiness, career orientation, and counseling support are able to improve the decision-making skills of students. However, most of the studies have been dedicated to career decision-making, psychological flexibility, and general counseling interventions. There were few studies which directly addressed value-centered counseling as a structured independent variable with dimensions like value clarification, self-awareness, responsibility, moral and social commitment, and future orientation. It is the main gap that the present research is supposed to address.

\section{Methodology}

\subsection{Research Approach and Design}

This research uses a quantitative research design with the objective of quantitatively analyzing the relationship between value-centered counseling and decision-making skills among university students. The research design used is descriptive-correlation survey that describes the levels of the two variables and analyzes how far the value-centered counseling influences the decision-making skills of the students. This design is most appropriate considering the use of questionnaire as the tool of data collection.

\subsection{Population and Sample}

The population is made up of students enrolled in Jordanian universities in the 2025/2026 academic year. The Ministry of Higher Education and Scientific Research reveals that Jordanian higher education has 480,009 students, out of which 287,320 are in public universities (60\%) while 192,689 are in private universities (40\%).

For this study, students from selected public and private universities in Jordan shall be targeted. The proposed public universities are The University of Jordan, Yarmouk University, Mutah University, and The Hashemite University; the proposed private universities are Middle East University and Amman Al-Ahliyya University. These institutions were chosen in order to achieve variation with respect to the type of university, academic environment, and student demographics. Student population data for each university and gender are available from the Department of Statistics based on the Statistical Yearbook of Jordan and Ministry of Higher Education and Scientific Research.

The sample will comprise of 384 students which is appropriate for a large population at 95\% confidence level and 5\% margin of error. Contemporary methodological literature recommends the use of population size, confidence level, margin of error, and sampling technique in determining sample size (Ahmed, 2024). In order to improve representativeness of the sample, allocation of the sample according to university type shall be proportional with about 230 students from public universities and 154 from private universities, representing 60\% and 40\%, respectively.

\subsection{Data Collection Instrument}

Data will be collected using a questionnaire containing three sections: demographic data, value-centered counseling, and decision-making skills. The value-centered counseling scale will cover the concepts of value clarification, self-awareness, personal responsibility, moral and social commitment, and future orientation. The decision-making skills scale will cover problem identification, information gathering, evaluation of alternatives, appropriate choice of decision, and confidence and responsibility in decision-making. A five-point Likert scale (strongly disagree to strongly agree) will be used.

\subsection{Validity and Reliability}

Content validity will be done by showing the questionnaire to experts in counseling, educational psychology, and research methodology with their views used in revising the questions before finalization of the questionnaire. Reliability will be assessed by use of Cranach's Alpha after conducting a pilot test with a few numbers of students that are not part of the final sample. Cranach's Alpha of 0.70 and above will be considered satisfactory.

\subsection{Statistical Methods}

Data analysis will be carried out using SPSS. Demographic variables will be described using descriptive statistics like frequencies and percentages. Value-centered counseling and decision-making skills levels will be described using means and standard deviation. Pearson correlation will be used to analyze the relationship between the two variables. Effect of value-centered counseling on decision-making skills will be tested through simple regression while multiple regressions will be used to test the effect of each aspect of value-centered counseling on decision-making skills.

\section{Results and Analysis}

\subsection{Demographic Characteristics of the Respondents}

This section presents the demographic profile of the study sample according to gender, age group, academic year, university type, and field of study. The total number of valid responses was 384, and no missing cases appeared in these demographic variables.


{\small
\begin{longtable}[]{@{}
  >{\raggedright\arraybackslash}p{(\columnwidth - 6\tabcolsep) * \real{0.2906}} >{\raggedright\arraybackslash}p{(\columnwidth - 6\tabcolsep) * \real{0.3584}} >{\raggedright\arraybackslash}p{(\columnwidth - 6\tabcolsep) * \real{0.1816}} >{\raggedright\arraybackslash}p{(\columnwidth - 6\tabcolsep) * \real{0.1693}}@{}}
\caption{Demographic Distribution of the Respondents}
\label{tab1}\\
\toprule\noalign{}
\endhead
\bottomrule\noalign{}
\endlastfoot
Variable & Category & Frequency & Percent \% \\
Gender & Male & 180 & 46.9 \\
 & Female & 204 & 53.1 \\
 & Total & 384 & 100.0 \\
Age Group & Less than 20 years & 75 & 19.5 \\
 & 20--22 years & 197 & 51.3 \\
 & 23--25 years & 83 & 21.6 \\
 & More than 25 years & 29 & 7.6 \\
 & Total & 384 & 100.0 \\
Academic Year & First year & 97 & 25.3 \\
 & Second year & 87 & 22.7 \\
 & Third year & 94 & 24.5 \\
 & Fourth year & 84 & 21.9 \\
 & Fifth year or above & 22 & 5.7 \\
 & Total & 384 & 100.0 \\
University Type & Public university & 230 & 59.9 \\
 & Private university & 154 & 40.1 \\
 & Total & 384 & 100.0 \\
Field of Study & Humanities and Social Sciences & 68 & 17.7 \\
 & Educational Sciences & 71 & 18.5 \\
 & Business and Economics & 81 & 21.1 \\
 & Sciences & 44 & 11.5 \\
 & Engineering and Technology & 65 & 16.9 \\
 & Medical and Health Sciences & 55 & 14.3 \\
 & Total & 384 & 100.0 \\
\end{longtable}
}

From the demographic results, there is relative parity in the distribution of gender with 53.1\% being female and 46.9\% male. In terms of age groups, most of the respondents were between the ages of 20-22 years (51.3\%), an expected trend in a sample of university students. The distribution of academic year was equal among the first four years while fifth year and above made up the smallest category (5.7\%). On the other hand, with respect to the university affiliation, students from public universities made up the larger category (59.9\%) compared to private universities (40.1\%).

\subsection{Reliability Analysis}

This section presents the reliability result of the study scale using Cronbach's Alpha. Reliability analysis was conducted to ensure that the items used in the questionnaire are internally consistent and suitable for further statistical analysis.

{\small
\begin{longtable}[]{@{}
  >{\raggedright\arraybackslash}p{(\columnwidth - 4\tabcolsep) * \real{0.2681}} >{\raggedright\arraybackslash}p{(\columnwidth - 4\tabcolsep) * \real{0.3830}} >{\raggedright\arraybackslash}p{(\columnwidth - 4\tabcolsep) * \real{0.3489}}@{}}
\caption{Reliability Statistics of the Study Scale}
\label{tab2}\\
\toprule\noalign{}
\endhead
\bottomrule\noalign{}
\endlastfoot
Scale & Cronbach's Alpha & Number of Items \\
Study Scale & 0.870 & 10 \\
\end{longtable}
}

The result of the reliability test is that Cronbach's Alpha is 0.870 for 10 questions. This shows very high internal consistency of the variables. The alpha is above the commonly used standard of 0.70, which indicates that the questions are reliable enough for use in the analysis.

\subsection{Descriptive Statistics of the Study Variables}

This section presents the descriptive statistics for the dimensions of value-centered counseling and decision-making skills. Means and standard deviations were used to identify the general level of respondents' agreement with each dimension.


{\small
\begin{longtable}[]{@{}
  >{\raggedright\arraybackslash}p{(\columnwidth - 12\tabcolsep) * \real{0.2021}} >{\raggedright\arraybackslash}p{(\columnwidth - 12\tabcolsep) * \real{0.1854}} >{\raggedright\arraybackslash}p{(\columnwidth - 12\tabcolsep) * \real{0.0589}} >{\raggedright\arraybackslash}p{(\columnwidth - 12\tabcolsep) * \real{0.1339}} >{\raggedright\arraybackslash}p{(\columnwidth - 12\tabcolsep) * \real{0.1428}} >{\raggedright\arraybackslash}p{(\columnwidth - 12\tabcolsep) * \real{0.0919}} >{\raggedright\arraybackslash}p{(\columnwidth - 12\tabcolsep) * \real{0.1850}}@{}}
\caption{Descriptive Statistics of Value-Centered Counseling and Decision-Making Skills Dimensions}
\label{tab3}\\
\toprule\noalign{}
\endhead
\bottomrule\noalign{}
\endlastfoot
Variable & Dimension & N & Minimum & Maximum & Mean & Std. Deviation \\
Value-Centered Counseling & Value Clarification & 384 & 1.25 & 5.00 & 3.5208 & 0.64928 \\
& Self-Awareness & 384 & 1.75 & 5.00 & 3.5508 & 0.65812 \\
& Personal Responsibility & 384 & 1.75 & 5.00 & 3.5150 & 0.66496 \\
& Moral and Social Commitment & 384 & 1.75 & 5.00 & 3.5449 & 0.64818 \\
& Future Orientation & 384 & 1.50 & 5.00 & 3.5521 & 0.65876 \\
Decision-Making Skills & Problem Identification & 384 & 1.50 & 5.00 & 3.5508 & 0.66183 \\
& Information Gathering & 384 & 1.75 & 5.00 & 3.5410 & 0.64592 \\
& Evaluation of Alternatives & 384 & 1.75 & 5.00 & 3.5443 & 0.66131 \\
& Choosing the Appropriate Decision & 384 & 1.50 & 5.00 & 3.5299 & 0.65791 \\
& Confidence and Responsibility & 384 & 1.50 & 5.00 & 3.5599 & 0.66892 \\
\end{longtable}
}

All the dimensions have been found to be having a mean score greater than the theoretical mean of 3.00, which represents the favorable degree of agreement of the respondents. In terms of value-based counseling, the highest mean has been recorded for Future Orientation (3.5521), whereas the lowest mean is recorded for Personal Responsibility (3.5150); but, at the same time, there is not much difference between dimensions. As far as decision-making skills are concerned, the highest mean is recorded for Confidence and Responsibility (3.5599), and Choosing the Appropriate Decision has the lowest mean (3.5299).

\subsection{Correlation Analysis}

The following is the output of the Pearson correlation coefficients for the value-centered counseling dimensions and decision-making skills dimensions. This test was done to find the nature of the relationship between the variables under study.

{\small
\begin{longtable}[]{@{}
  >{\raggedright\arraybackslash}p{(\columnwidth - 10\tabcolsep) * \real{0.1753}} >{\raggedright\arraybackslash}p{(\columnwidth - 10\tabcolsep) * \real{0.1753}} >{\raggedright\arraybackslash}p{(\columnwidth - 10\tabcolsep) * \real{0.1565}} >{\raggedright\arraybackslash}p{(\columnwidth - 10\tabcolsep) * \real{0.1659}} >{\raggedright\arraybackslash}p{(\columnwidth - 10\tabcolsep) * \real{0.1515}} >{\raggedright\arraybackslash}p{(\columnwidth - 10\tabcolsep) * \real{0.1754}}@{}}
\caption{Pearson Correlations between Value-Centered Counseling and Decision-Making Skills Dimensions}
\label{tab4}\\
\toprule\noalign{}
\endhead
\bottomrule\noalign{}
\endlastfoot
Value-Centered Counseling Dimensions & Problem Identification & Information Gathering & Evaluation of Alternatives & Choosing the Appropriate Decision & Confidence and Responsibility \\
Value Clarification & .426** & .383** & .444** & .424** & .434** \\
Self-Awareness & .433** & .329** & .397** & .430** & .385** \\
Personal Responsibility & .346** & .355** & .358** & .372** & .358** \\
Moral and Social Commitment & .286** & .319** & .337** & .373** & .420** \\
Future Orientation & .385** & .374** & .411** & .430** & .434** \\
\end{longtable}
}

Note. N = 384. Correlation is significant at the 0.01 level.

It should be noted that all the relationships found between the dimensions of value-centered counseling and decision-making skills are positive and statistically significant (0.01). Thus, a greater level of value clarification, self-awareness, responsibility, moral and social commitment and future orientation are connected with higher decision-making skills in students. Among the strongest connections are those found between Value Clarification and Evaluation of Alternatives (r = .444), Self-Awareness and Problem Identification (r = .433) and Future Orientation with Confidence and Responsibility (r = .434). In general, the relationships are rather moderate but positive, which indicates the presence of an association between value-centered counseling and decision-making skills.

\subsection{Testing Hypotheses by Means of Multiple Regression}

This part provides the results of hypothesis testing about the influence of value-centered counseling dimensions on decision-making skills dimensions. The hypothesis is considered to be proven if the corresponding significance level is less than 0.05.

{\small
\begin{longtable}[]{@{}
  >{\raggedright\arraybackslash}p{(\columnwidth - 10\tabcolsep) * \real{0.1832}} >{\raggedright\arraybackslash}p{(\columnwidth - 10\tabcolsep) * \real{0.2908}} >{\raggedright\arraybackslash}p{(\columnwidth - 10\tabcolsep) * \real{0.2093}} >{\raggedright\arraybackslash}p{(\columnwidth - 10\tabcolsep) * \real{0.0928}} >{\raggedright\arraybackslash}p{(\columnwidth - 10\tabcolsep) * \real{0.0735}} >{\raggedright\arraybackslash}p{(\columnwidth - 10\tabcolsep) * \real{0.1504}}@{}}
\caption{Results of Hypotheses Testing for the Effect of Value-Centered Counseling Dimensions on Decision-Making Skills}
\label{tab5}\\
\toprule\noalign{}
\endhead
\bottomrule\noalign{}
\endlastfoot
Hypothesis & Relationship Tested & Dependent Variable & Beta & Sig. & Decision \\
H2a-1 & Value Clarification → Problem Identification & Problem Identification & .247 & .000 & Supported \\
H2a-2 & Value Clarification → Information Gathering & Information Gathering & .210 & .000 & Supported \\
H2a-3 & Value Clarification → Evaluation of Alternatives & Evaluation of Alternatives & .262 & .000 & Supported \\
H2a-4 & Value Clarification → Choosing the Appropriate Decision & Choosing the Appropriate Decision & .213 & .000 & Supported \\
H2a-5 & Value Clarification → Confidence and Responsibility & Confidence and Responsibility & .228 & .000 & Supported \\
H2b-1 & Self-Awareness → Problem Identification & Problem Identification & .230 & .000 & Supported \\
H2b-2 & Self-Awareness → Information Gathering & Information Gathering & .089 & .082 & Not Supported \\
H2b-3 & Self-Awareness → Evaluation of Alternatives & Evaluation of Alternatives & .152 & .002 & Supported \\
H2b-4 & Self-Awareness → Choosing the Appropriate Decision & Choosing the Appropriate Decision & .184 & .000 & Supported \\
H2b-5 & Self-Awareness → Confidence and Responsibility & Confidence and Responsibility & .108 & .025 & Supported \\
H2c-1 & Personal Responsibility → Problem Identification & Problem Identification & .117 & .015 & Supported \\
H2c-2 & Personal Responsibility → Information Gathering & Information Gathering & .159 & .001 & Supported \\
H2c-3 & Personal Responsibility → Evaluation of Alternatives & Evaluation of Alternatives & .123 & .010 & Supported \\
H2c-4 & Personal Responsibility → Choosing the Appropriate Decision & Choosing the Appropriate Decision & .127 & .007 & Supported \\
H2c-5 & Personal Responsibility → Confidence and Responsibility & Confidence and Responsibility & .110 & .018 & Supported \\
H2d-1 & Moral and Social Commitment → Problem Identification & Problem Identification & .020 & .675 & Not Supported \\
H2d-2 & Moral and Social Commitment → Information Gathering & Information Gathering & .105 & .037 & Supported \\
H2d-3 & Moral and Social Commitment → Evaluation of Alternatives & Evaluation of Alternatives & .085 & .075 & Not Supported \\
H2d-4 & Moral and Social Commitment → Choosing the Appropriate Decision & Choosing the Appropriate Decision & .121 & .010 & Supported \\
H2d-5 & Moral and Social Commitment → Confidence and Responsibility & Confidence and Responsibility & .195 & .000 & Supported \\
H2e-1 & Future Orientation → Problem Identification & Problem Identification & .184 & .000 & Supported \\
H2e-2 & Future Orientation → Information Gathering & Information Gathering & .192 & .000 & Supported \\
H2e-3 & Future Orientation → Evaluation of Alternatives & Evaluation of Alternatives & .211 & .000 & Supported \\
H2e-4 & Future Orientation → Choosing the Appropriate Decision & Choosing the Appropriate Decision & .219 & .000 & Supported \\
H2e-5 & Future Orientation → Confidence and Responsibility & Confidence and Responsibility & .229 & .000 & Supported \\
\end{longtable}
}

{\small
\begin{longtable}[]{@{}
  >{\raggedright\arraybackslash}p{(\columnwidth - 6\tabcolsep) * \real{0.5366}} >{\raggedright\arraybackslash}p{(\columnwidth - 6\tabcolsep) * \real{0.1373}} >{\raggedright\arraybackslash}p{(\columnwidth - 6\tabcolsep) * \real{0.0754}} >{\raggedright\arraybackslash}p{(\columnwidth - 6\tabcolsep) * \real{0.2507}}@{}}
\caption{Model-Level Results}
\label{tab6}\\
\toprule\noalign{}
\endhead
\bottomrule\noalign{}
\endlastfoot
Dependent Variable & F Value & Sig. & Model Decision \\
Problem Identification & 35.805 & .000 & Significant \\
Information Gathering & 28.165 & .000 & Significant \\
Evaluation of Alternatives & 38.239 & .000 & Significant \\
Choosing the Appropriate Decision & 41.892 & .000 & Significant \\
Confidence and Responsibility & 42.765 & .000 & Significant \\
\end{longtable}
}

The results reveal that all regression models attained statistical significance at p \textless{} .001, which indicates that all dimensions of value-centered counseling have a significant impact on all dimensions of decision-making skills. Hence, the overall hypothesis of the current research is confirmed.

As for specific hypotheses, the Value Clarification hypothesis was supported in all five relationships and showed significant prediction of all dimensions of decision-making skills. Future Orientation was confirmed in all five relationships and showed a significant influence on all decision-making dimensions. Also, the Personal Responsibility hypothesis was supported in all five relationships with all dependent variables.

Nevertheless, several hypotheses were partially supported. First, Self-Awareness was significant in four relationships but did not significantly predict Information Gathering (p = .082, above the level of .05). So, the respective hypothesis is not confirmed. Second, Moral and Social Commitment was partially supported in that it significantly predicted three dimensions, including Information Gathering, Choosing the Appropriate Decision, and Confidence and Responsibility but did not significantly predict Problem Identification and Evaluation of Alternatives (p = .675 and p = .075). Consequently, these two hypotheses are not confirmed.

To summarize, the research findings generally confirm the formulated hypotheses, although the negative results indicate that not all dimensions of value-centered counseling affect each dimension of decision-making skills equally.

\section{Discussion}

According to the research findings, value-centered counseling has been positively perceived among Jordanian university students with all its dimensions having means above the theoretical midpoint. In addition, the Future Orientation dimension had the highest mean value, which implies that students perceive counseling activities as beneficial mostly when they help to link current decisions to the future academic and career goals. Thus, the findings support the theoretical assumption that values are the criteria guiding behaviors and making actions more meaningful \cite{ref11,ref14}. The findings also support the notion that counseling must go beyond problem-solving and promote awareness, responsibility, and personal direction.

Moreover, the research shows that decision-making skills were positively perceived by Jordanian university students, and all five dimensions of decision-making skills had means above the midpoint, but Confidence and Responsibility had the highest means whereas choosing the Appropriate Decision had the lowest mean. This implies that even though students have adequate decision-making skills, it might be hard for them to make decisions because of the uncertainty and sense of responsibility in their choices. The finding is consistent with the arguments of Gati and Kulcsár \cite{ref15} that efficient decision-making consists of identification of alternatives, their evaluation, and selection of the most suitable one. Besides, the research findings resonate with Le et al. \cite{ref35} who claim that decision-making is one of the major life skills needed to cope with university challenges.

The correlation analysis performed has shown positive and significant relationships between all dimensions of value-centered counseling and decision-making skills, which implies that the first hypothesis of the research has been confirmed \cite{ref40}. Namely, the higher the value clarification, self-awareness, responsibility, moral and social commitment, and future orientation were, the higher decision-making skills students had. The finding is consistent with the findings of \cite{ref3} according to who values clarification helps to prepare oneself and be confident in making decisions. \cite{ref36} Link human values with decision-making styles.

Moreover, the performed regression analysis has confirmed that value-centered counseling significantly predicts decision-making skills because all models were significant at p \textless{} .001, and thus the main hypothesis about the statistically significant effect of value-centered counseling on decision-making skills has been confirmed. The result is consistent with the literature proving that values-based counseling and committed action can positively affect reflection, psychological flexibility, and responsible behavior \cite{ref23,ref25}. Furthermore, the research finding coincides with the result of \cite{ref34} who finds that counseling programs can improve decision-making skills among Jordanian university students.

As for the relationships between individual dimensions of both constructs, the Value Clarification, Personal Responsibility, and Future Orientation were identified as the most consistent and strong predictors because they significantly affected all five dimensions of decision-making skills. This implies that students make better decisions when they know their values, take responsibility for decisions, and foresee future consequences. The finding is confirmed by \cite{ref27} who reports that career value clarification improves career maturity and decision-making of individuals. Moreover, the research finding is confirmed by the work of \cite{ref28} who find out that career values and self-efficacy improve the quality of students' career decisions.

However, some of the research findings remain only partially confirmed. First, the Self-Awareness did not significantly predict Information Gathering, which means that knowing oneself does not always imply proactive seeking of information from outside sources \cite{ref39}. Second, Moral and Social Commitment did not significantly predict Problem Identification or Evaluation of Alternatives, which implies that moral and social values have a higher effect on selection and accepting of responsibility rather than on analytical stages of problem definition and alternatives comparison \cite{ref38}. The conclusion is consistent with \cite{ref18} who state that decision-making difficulties can emerge due to the lack of information and readiness rather than values.

To sum up, the research proves that value-centered counseling plays an important role in the improvement of decision-making skills among Jordanian university students. The research findings prove that decision-making is not a purely rational or technical process but a process based on values and psychology, and it is characterized by the clarity of values, awareness, responsibility, confidence, and future orientation. Consequently, Jordanian universities should develop counseling services that will help students clarify their values, develop responsibility, evaluate alternatives, and make decisions according to their academic and personal goals.

\section{Conclusion}

The research has concluded that value-centered counseling has a significant impact on the improvement of decision-making skills among Jordanian university students. The research results prove the positive levels of value-centered counseling and decision-making skills, and significant correlations have been found across all dimensions. Regression analysis has proved that value-centered counseling significantly predicts students' ability to identify problems, gather information, evaluate alternatives, make appropriate decisions, and take responsibility for decisions.

Three dimensions, including Value Clarification, Personal Responsibility, and Future Orientation, have proven to be the most consistent predictors of decision-making skills which imply that students make better decisions when they know their values, take responsibility, and see future implications of current decisions.

\textbf{Recommendations}

The study has recommended to Jordanian universities that they should develop counseling services focused on values clarification, self-awareness, responsibility, and future plans. The university counseling centers should implement training programs that will help students to align their academic and personal decisions with values and future plans.

Besides, it is recommended to integrate decision-making skills into student development programs, particularly those concerned with the development of competencies in problem identification, information gathering, evaluation of alternatives, and responsible decision making. Counseling should not only guide students but also equip them with the ability to think independently and make value-based decisions.

In addition, attention should be paid to students who hesitate in decision-making, have low confidence in making decisions, or rely on other people's help in this process. Workshops, group counseling sessions, and academic advising can help to increase students' confidence and responsibility.

Finally, the future research should involve using experimental design or qualitative interviews to get deeper insight into the effect of value-centered counseling on students' decision-making in practice

\textbf{Conflicts of Interest:}

The authors declare no conflict of interest

\textbf{Funding:}

This study was funded by Deanship of Scientific Research at Jadara University.

\textbf{Open Access:} The researchers adhered to strict ethical principles, ensuring data confidentiality and protecting the rights and safety of participants throughout the study, so that this research can be a reference for all.

\textbf{Data Availability Statement:}

The data supporting the findings of this study are available from the corresponding author upon reasonable request; all datasets were collected and stored in compliance with institutional research ethics and privacy guidelines

\textbf{ACKNOWLEDGEMENTS:}

The authors are grateful to the Deanship of Scientific Research at Jadara University for providing financial support for this publication"

\emph{\textbf{Conflicts of Interest Statement}}

\emph{The authors certify that they have NO affiliations with or involvement in any organization or entity with any financial interest (such as honoraria; educational grants; participation in speakers' bureaus; membership, employment, consultancies, stock ownership, or other equity interest; and expert testimony or patent-licensing arrangements), or non-financial interest (such as personal or professional relationships, affiliations, knowledge or beliefs) in the subject matter or materials discussed in this manuscript.}

\begin{thebibliography}{99}

\bibitem{ref1} Lee, S., Jung, J., Baek, S., \& Lee, S. (2022). The relationship between career decision-making self-efficacy, career preparation behaviour and career decision difficulties among South Korean college students. \emph{Sustainability, 14}(21), 14384. doi: 10.3390/su142114384

\bibitem{ref2} Moore, B., \& Gregory, R. (2024). Decision making as a pedagogy for social emotional learning. \emph{Social and Emotional Learning: Research, Practice, and Policy, 3}, 100034. DOI: 10.1016/j.sel.2024.100034

\bibitem{ref3} Witteman, H. O., Ndjaboue, R., Vaisson, G., Dansokho, S. C., Arnold, B., Bridges, J. F. P., Comeau, S., Fagerlin, A., Gavaruzzi, T., Marcoux, M., Pieterse, A., Pignone, M., Provencher, T., Racine, C., Regier, D., Rochefort-Brihay, C., Thokala, P., Weernink, M., White, D. B., Wills, C. E., \& Jansen, J. (2021). Clarifying values: An updated and expanded systematic review and meta-analysis. \emph{Medical Decision Making, 41}(7), 801--820. DOI: 10.1177/0272989X211037946

\bibitem{ref4} Gunawan, I. M. S., Wibowo, M. E., Purwanto, E., \&Sunawan, S. (2019). Group counseling of values clarification to increase middle school students' empathy. \emph{PsicologíaEducativa, 25}(2), 169--174. DOI: 10.5093/psed2019a5

\bibitem{ref5} Sagone, E., \& Indiana, M. L. (2021). Are decision-making styles, locus of control, and average grades in exams correlated with procrastination in university students? \emph{Education Sciences, 11}(6), 300. DOI: 10.3390/educsci11060300

\bibitem{ref6} Darawsheh , N.(2023). The Impact of Cyber Bullying on the Psychological Well-being of University Students: A Study in Jordanian Universities, \emph{Information Sciences Letters}, 8 (2), 2757-2768

\bibitem{ref7} Sutiman, Sofyan, H., Soenarto, Mutohhari, F., \&Nurtanto, M. (2022). Students' career decision-making during online learning: The mediating roles of self-efficacy in vocational education. \emph{European Journal of Educational Research,} 11(3), 1669--1682. DOI: 10.12973/eu-jer.11.3.1669

\bibitem{ref8} Lipshits-Braziler, Y., Arieli, S., \& Daniel, E. (2025). Personal values and career-related preferences among young adults. \emph{Journal of Personality, 93}(2), 378--393. DOI: 10.1111/jopy.12935

\bibitem{ref9} Darawsheh, N , Alrashdan, H , Al-Khataybeh, M , Al Ajmi, H , almheiri, A , Shaglof,A , Aldarmaki, I , BaniYounes, Z and Aidi,R .(2025). The Degree of Impact of Organizational Intelligence Obstacles in the Digital Age, Journal of Statistics Applications \& Probability, \emph{An International Journal}, 14(6), 985-995.

\bibitem{ref10} \href{https://www.scopus.com/authid/detail.uri?authorId=57299687800}{ALAwAmRAh, A} , \href{https://www.scopus.com/authid/detail.uri?authorId=59210073600}{Ajmi, H}, \href{https://www.scopus.com/authid/detail.uri?authorId=59171231500}{Darawsheh, N} , \href{https://www.scopus.com/authid/detail.uri?authorId=58808922400}{Darawsheh, O}, \href{https://www.scopus.com/authid/detail.uri?authorId=59210786400}{Urabi, O.}(2024).\emph{Analyzing Student's Perceptions about the Values of Digital Citizenship,} Journal of Statistics Applications and Probability\emph{,\textbf{~}}2024, 13(4), pp. 1279--1288\emph{\textbf{.}}

\bibitem{ref11} Berkout, O. V. (2022). Working with values: An overview of approaches and considerations in implementation. \emph{Behavior Analysis in Practice, 15}(1), 104--114. \href{https://doi.org/10.1007/s40617-021-00589-1}{\ul{https://doi.org/10.1007/s40617-021-00589-1}}

\bibitem{ref12} Klussman, K., Curtin, N., Langer, J., \& Nichols, A. L. (2022). The importance of awareness, acceptance, and alignment with the self: A framework for understanding self-connection. \emph{Europe's Journal of Psychology, 18}(1), Article e3707. \href{https://doi.org/10.5964/ejop.3707}{\ul{https://doi.org/10.5964/ejop.3707}}

\bibitem{ref13} Alrashdan, H; Darawsheh, N; Aabed, S; Bsoul, M; Ganayem, S; and Mhameed, M. (2023). Analyzing University Administration's Impact on Faculty Digital Self-Learning, Journal of Statistics Applications and Probability, 12, 1591--1601

\bibitem{ref14} Ryan, R. M., \& Deci, E. L. (2020). Intrinsic and extrinsic motivation from a self-determination theory perspective: Definitions, theory, practices, and future directions. \emph{Contemporary Educational Psychology, 61}, 101860. \href{https://doi.org/10.1016/j.cedpsych.2020.101860}{\ul{https://doi.org/10.1016/j.cedpsych.2020.101860}}

\bibitem{ref15} Gati, I., \& Kulcsár, V. (2021). Making better career decisions: From challenges to opportunities. \emph{Journal of Vocational Behavior, 126}, 103545. \href{https://doi.org/10.1016/j.jvb.2021.103545}{\ul{https://doi.org/10.1016/j.jvb.2021.103545}}

\bibitem{ref16} Alrashdan, H., Darawsheh, N., Almheiri, A., Aljanabi, Y., Al-Khataybeh, M., \& Al-Quran, N. (2025). The Effect of Administrative Bullying on the Quality of Education. \emph{Educational Process: International} \emph{Journal, 18}, e2025515.

\bibitem{ref17} Darawsheh, N. (2018). The Reality of The Challenges Faced By Graduate Students In The Faculties Of Educational Sciences In Jordanian Universities, \emph{JIRSEA,} 16(2),132-140,(2018).

\bibitem{ref18} Kulcsár, V., Dobrean, A., \& Gati, I. (2020). Challenges and difficulties in career decision making: Their causes, and their effects on the process and the decision. \emph{Journal of Vocational Behavior, 116}, 103346. \href{https://doi.org/10.1016/j.jvb.2019.103346}{\ul{https://doi.org/10.1016/j.jvb.2019.103346}}

\bibitem{ref19} Azhenov, A., Kudysheva, A. K., Fominykh, N. F., \&Tulekova, G. T. (2023). Career decision-making readiness among students' in the system of higher education: Career course intervention. \emph{Frontiers in Education, 8}, 1097993. \href{https://doi.org/10.3389/feduc.2023.1097993}{\ul{https://doi.org/10.3389/feduc.2023.1097993}}

\bibitem{ref20} Hatamala, H ,Darawsheh, N. (2019). The Challenges of the Application of the Productive University's Philosophy In Jordanian Universities and Ways of Developing Them from The Perspective of Academic Leaders,\emph{\href{https://www.scopus.com/authid/detail.uri?authorId=57205221009\#disabled}{Journal of Institutional Research South East Asia}\textbf{, }}17(1), pp. 93--117,(2019).

\bibitem{ref21} Darawsheh, N; Al-Khataybeh, M; Alrashdan, H; Drawsheh, S; Abu Grarh, A; and abu dogush, M .(2024) "Coexistence skills among mothers of children with multiple disabilities in the northern region of Jordan (IRBID)," \emph{\textbf{Journal of}} \emph{\textbf{Statistics Applications \& Probability}}: 13( 2), 665-680

\bibitem{ref22} Bayly, B. L., \& Bumpus, M. F. (2020). Predictors and implications of values clarity in first-year college students. \emph{College Student Journal, 53}(4), 397--404.

\bibitem{ref23} Viskovich, S., Pakenham, K. I., \& Fowler, J. A. (2021). A mixed-methods evaluation of experiential intervention exercises for values and committed action from an Acceptance and Commitment Therapy mental health promotion program for university students. \emph{Journal of Contextual Behavioral Science, 22}, 108--118. doi: 10.1016/j.jcbs.2021.10.001

\bibitem{ref24} Larsson, A., Hartley, S., \& McHugh, L. (2022). A randomised controlled trial of brief web-based acceptance and commitment therapy on the general mental health, depression, anxiety and stress of college students. \emph{Journal of Contextual Behavioral Science, 24}, 10--17. doi: 10.1016/j.jcbs.2022.02.005

\bibitem{ref25} Hsu, T., Adamowicz, J. L., \& Thomas, E. B. K. (2023). The effect of acceptance and commitment therapy on the psychological flexibility and inflexibility of undergraduate students: A systematic review and three-level meta-analysis. \emph{Journal of Contextual Behavioral Science, 30}, 169--180. doi: 10.1016/j.jcbs.2023.10.006

\bibitem{ref26} De Gagne, J. C. (2023). Values clarification exercises to prepare nursing students for artificial intelligence integration. \emph{International Journal of Environmental Research and Public Health, 20}(14), 6409. doi: 10.3390/ijerph20146409

\bibitem{ref27} Heo, J. (2023). Effect of a career value clarification program on career consciousness maturity and career decision-making of female college students. \emph{Journal of The Korea Society of Computer and Information, 28}(11), 191--199. doi: 10.9708/jksci.2023.28.11.191

\bibitem{ref28} Wang, X.-H., Wang, H.-P., \& WenYa, L. (2023). Improving the quality of career decision-making of students in Chinese higher vocational colleges. \emph{SAGE Open, 13}(2), 1--14. doi: 10.1177/21582440231180105

\bibitem{ref29} Liu, X., Zhang, X., Dang, Y., \& Gao, W. (2023). Career education skills and career adaptability among college students in China: The mediating role of career decision-making self-efficacy. \emph{Behavioral Sciences, 13}(9), 780. doi: 10.3390/bs13090780

\bibitem{ref30} Salim, R. M. A., Istiasih, M. R., Rumalutur, N. A., \& Situmorang, D. D. B. (2023). The role of career decision self-efficacy as a mediator of peer support on students' career adaptability. \emph{Heliyon, 9}(4), e14911. doi: 10.1016/j.heliyon.2023.e14911

\bibitem{ref31} Zhou, A., Liu, J., Xu, C., \& Jobe, M. C. (2024). Effect of social support on career decision-making difficulties: The chain mediating roles of psychological capital and career decision-making self-efficacy. \emph{Behavioral Sciences, 14}(4), 318. doi: 10.3390/bs14040318

\bibitem{ref32} Pham, M., Lam, B. Q., \& Bui, A. T. N. (2024). Career exploration and its influence on the relationship between self-efficacy and career choice: The moderating role of social support. \emph{Heliyon, 10}(11), e31808. doi: 10.1016/j.heliyon.2024.e31808

\bibitem{ref33} Milledzi, E. Y., Asamani, L., Opoku Mensah, A., \&Naamwanuru, V. (2024). Career decision self-efficacy mediates the relationship between general self-efficacy and career decision-making of undergraduate students. \emph{Canadian Journal of Career Development, 23}(2), 169--181. doi: 10.53379/cjcd.2024.389

\bibitem{ref34} Daheri, S. (2022). The effectiveness of a group counseling program on developing decision making skill of educational sciences students at counseling and mental health department. \emph{The Online Journal of New Horizons in Education, 12}(3), 204--210.

\bibitem{ref35} Le, H. T. T., Pham, L. T., \& Vu, H. T. T. (2022). Student decision-making processes as evaluated by students, administrators, and lecturers. \emph{International Journal of Education and Practice, 10}(4), 371--380. DOI: 10.18488/61.v10i4.3221

\bibitem{ref36} Páez Gallego, J., De-Juanas Oliva, Á., García-Castilla, F. J., \& Muelas, Á. (2020). Relationship between basic human values and decision-making styles in adolescents. \emph{International Journal of Environmental Research and Public Health, 17}(22), 8315. \href{https://doi.org/10.3390/ijerph17228315}{\ul{https://doi.org/10.3390/ijerph17228315}}

\bibitem{ref37} Christodoulou, V., Flaxman, P. E., \& Lloyd, J. (2021). Acceptance and commitment therapy in group format for college students. \emph{Journal of College Counseling, 24}(3), 210--223. \href{https://doi.org/10.1002/jocc.12192}{\ul{https://doi.org/10.1002/jocc.12192}}

\bibitem{ref38} Ayre, J., Jenkins, H., Kumarage, R., McCaffery, K. J., Maher, C. G., \& Hancock, M. J. (2025). Exploring values clarification and health-literate design in patient decision aids: A qualitative interview study. \emph{Medical Decision Making, 45}(5), 510--521. doi: 10.1177/0272989X251334356

\bibitem{ref39} Ahmed, S. K. (2024). How to choose a sampling technique and determine sample size for research: A simplified guide for researchers. \emph{Oral Oncology Reports, 12}, 100662. \url{https://doi.org/10.1016/j.oor.2024.100662}

\bibitem{ref40} Krejcie, R. V., \& Morgan, D. W. (1970). Determining sample size for research activities. \emph{Educational and Psychological Measurement, 30}(3), 607--610. \url{https://doi.org/10.1177/001316447003000308}

\end{thebibliography}

\end{document}
